\documentclass[11pt,a4paper]{article}
\usepackage{kotex}
\usepackage{amsmath,amssymb,booktabs,geometry}
\geometry{margin=25mm}
\title{I01 선검증 2: 검출 효과와 표지 분포를 바꾼 반증 실험}
\author{RELATIONAL ATLAS}
\date{2026년 10월 4일}
\begin{document}
\maketitle

\section{결론부터}
발전연구의 주제는 선택하지 않았다. 선검증 1의 조건부 패턴은
서로 다른 에너지 스펙트럼과 비투영 검출 효과에서도 유지됐지만,
기록된 표지의 \emph{실제 표본 분포}가 균일하다는 전제를 버리면
``표지 삭제 후 공통 위상 정보가 정확히 0''이라는 결론이 깨졌다.
반대로 $P=0$ 상태와 전하 twirl 검출기는 표지 분포와 무관하게
공통 위상 불변이었다. 이는 관측 rank가 상태, 검출 효과, 표지 설계의
교차 결과라는 선검증이며 공간 차원 선택의 증거가 아니다.

\section{사전 고정 설계와 식별 조건}
실행 전에 \texttt{prevalidation\_protocol2\_2026-10-04.tex}에
$n=3,4,5$, $\beta=0.45,1.20$, 세 검출 효과, 세 시계 스펙트럼,
$\tau=0,\pi/2$, 두 표지 분포, 비제약/$P=0$ 상태를 고정했다.
총 432행이며 독립된 432회 통계 실험이 아니라 결정론적 매개변수 격자다.
표지는 $L=2n+1$개, 검출 바이어스는 1.25 rad다.
원래 에너지 $E_0=\sum_i k_i^2$ 외에
\begin{equation}
E_{\rm nonlin}=E_0+0.37E_0^2/n,\qquad
E_{\rm charge}=E_0+0.41P^2/n
\end{equation}
을 시계 위상에 넣었다. 상태 가중치 $w_k\propto e^{-\beta E_0}$는
고정했으므로 상태 QFI 검사는 새 스펙트럼에 대해 독립된 새로운 상태
시험이 아니다.

복귀 기준 상태의 폭을 $\beta_r=0.15$로 바꾼 단일 프로젝터
$M_{\rm mis}$, $\beta_r=0.15,1.80$ 프로젝터의 0.7/0.3 혼합
$M_{\rm mix}$, 그리고 $M_{\rm mix}$의 총전하 twirl
$M_{\rm twirl}$을 비교했다. 혼합은 하나의 프로젝터가 아닌 유효
이진 POVM 효과다. 표지 분포는 균일 $1/L$ 또는
\begin{equation}
a_t=\frac{1+0.35\cos(2\pi t/L)}{L}
\end{equation}
이다. 두 분포 모두 표지값 자체는 같은 $2\pi t/L$을 쓴다.
따라서 차이는 표지를 얼마나 자주 사용했는가와 그것을 지운 후
어떤 혼합을 얻는가에만 있다.

전하 $P$의 차이가 만드는 Fourier 조화는
$|P-P'|\le2n$이다. 균일한 $L=2n+1$ 표지 삭제는 비영 조화를
모두 소거한다. 이 주장은 $E(k)$의 형태에는 의존하지 않는다.
그러나 위의 비균일 가중치는 $m\equiv\pm1\pmod L$ 조화를 남긴다.
현재 cutoff에서는 $m=\pm1$뿐 아니라 $m=\pm2n$도 여기에 해당한다.
$P=0$ 상태는 처음부터 공통 위상 조화가 없고,
$M_{\rm twirl}$은 전하 사이 coherence를 검출하지 않으므로
어느 표지 분포에서도 공통 위상에 불변이다.

\section{정확 격자와 Qiskit 교차검증}
\begin{center}
\begin{tabular}{@{}lrr@{}}
\toprule
검사 & 대상 행 & 실패/양성 행 \\
\midrule
사전 보장된 삭제 후 불변성과 rank $n-1$ & 360 & 실패 0 \\
표지 보존의 예상 rank & 432 & 실패 0 \\
상태 QFI의 비제약/$P=0$ rank $n/(n-1)$ & 432 & 실패 0 \\
비균일 표지, 비제약 coherent 효과의 잔여 반응 & 72 & 양성 72 \\
비균일 표지, 비제약 coherent 효과의 삭제 rank $n$ & 72 & 양성 72 \\
\bottomrule
\end{tabular}
\end{center}
공통 위상을 $0\to0.23$ rad 옮긴 뒤 최대 확률 변화의
$10^{-10}$ 판정선에서 보장된 불변성은 모두 유지됐다.
비균일/coherent 72행의 잔여 확률 변화는
$3.25\times10^{-7}$--$2.67\times10^{-3}$였으므로,
정확히 양수여도 샷 예산에 따라 숨을 수 있다. Fisher rank는 최대
고유값 대비 $10^{-8}$ 기준으로 보고했으며 공통방향 정보와
이동 전후 확률도 원자료에 보존했다. 해석적 확률 미분의 중앙차분
대조 최대 오차는 $1.71\times10^{-11}$이다.

선택한 $n=3$ 조건 네 개의 별도 Qiskit Statevector 회로에서는
시계 에너지 위상을 직접 대각 게이트로, 검출 효과를 역준비 회로와
혼합으로 구현했다. $M_{\rm twirl}$은 $L$개 공통 위상의 회로
확률 평균으로 계산해 전하-섹터 식과 교차 확인했다.
수식과 Statevector 확률 최대 차이는 $2.43\times10^{-15}$,
세 조건의 Aer density-matrix 최대 차이는 $3.05\times10^{-15}$였다.
비균일 표지의 회로 확률을 실제 $a_t$로 다시 합친 대조의
표지별 최대 차이는 $4.72\times10^{-15}$였다.
이는 구현 경로의 대조이며 유한 샷 검정력이나 QPU 성능을
보여주는 실험이 아니다.

\section{검증 상태와 남은 게이트}
현재 검증된 것은 cutoff 1의 특정 상태족에서 전하 Fourier 조화,
검출 효과의 coherence, 표지 분포가 관측 Fisher rank에 미치는
구조다. 확립된 발전연구의 주제는 아직 없다. 특히 이번 격자는
이미 알려진 표지 분포를 정확히 사용했다. 분포의 추정 오차,
기준 위상 자체의 오차, 드리프트 법칙의 오가정, cutoff 변경,
유한 샷에서 작은 잔여 신호의 검정력은 여전히 검증 게이트다.
공간방향의 실제 개수와 관측 rank를 분리한다는 기존 판정도 유지한다.
QPU는 사용하지 않았다.

\paragraph{재현 자료.}
사전 프로토콜: \texttt{prevalidation\_protocol2\_2026-10-04.tex}.
정확 격자: \texttt{../src/relational\_atlas\_i01/holdout/qiskit\_rank\_holdout2\_grid.py},
\texttt{../results/qiskit\_rank\_holdout2\_grid.json}.
독립 회로: \texttt{../src/relational\_atlas\_i01/holdout/qiskit\_rank\_holdout2\_circuit.py},
\texttt{../results/qiskit\_rank\_holdout2\_circuit.json}.
검증 테스트: \texttt{../tests/test\_qiskit\_rank\_holdout2\_grid.py},
\texttt{../tests/test\_qiskit\_rank\_holdout2\_circuit.py}.
\end{document}
